\chapter{Objetivos, Alcance y Antecedentes}
\label{cap:objetivos_alcance_antecedentes}

En este capítulo se establece la gobernanza científica, metodológica y contextual del Trabajo de Fin de Máster. En primer lugar, se formulan el objetivo general y los objetivos específicos del proyecto con un enfoque teleológico orientado a preguntas de investigación e hipótesis verificables, dando respuesta directa a las directrices de evaluación académica. Seguidamente, se delimitan rigurosamente las fronteras de la investigación, distinguiendo las inclusiones operativas de las exclusiones impuestas por el estado actual de la tecnología NISQ (\textit{Noisy Intermediate-Scale Quantum}). Por último, se revisa el estado del arte en optimización clásica y variacional cuántica, identificando las limitaciones fundamentales de los enfoques convencionales y fundamentando la necesidad de la arquitectura con preservación de simetría que se desarrollará formalmente en el Capítulo~\ref{cap:marco_teorico}.

% ==============================================================================
% 2.1. OBJETIVOS DEL PROYECTO
% ==============================================================================
\section{Objetivos del Proyecto}
\label{sec:objetivos_proyecto}

\subsection{Objetivo General}
\label{subsec:objetivo_general}

El presente Trabajo de Fin de Máster tiene como finalidad determinar en qué medida un algoritmo cuántico variacional con preservación de simetría estructural (XY-QAOA regularizado) permite resolver el problema de selección de carteras financieras bajo una restricción de cardinalidad exacta $K$ de forma más robusta ---en términos de factibilidad de las soluciones y calidad financiera de las mismas--- que el paradigma estándar de penalizaciones cuadráticas (QAOA), delimitando honestamente el rango de escalabilidad en el que dicha ventaja resulta observable bajo simulación clásica de vectores de estado, y contrastando sus resultados frente a solucionadores de referencia de uso industrial (Gurobi Optimizer) y metaheurísticas clásicas (Simulated Annealing).

\subsection{Objetivos Específicos}
\label{subsec:objetivos_especificos}

Para materializar y evaluar de manera objetiva la finalidad formulada en el objetivo general, el trabajo se estructura en torno a cuatro objetivos específicos (OE1--OE4), diseñados como ejes de investigación cuantitativos y verificables:

\begin{itemize}
    \item \textbf{OE1. Verificar la equivalencia matemática del mapeo Markowitz $\rightarrow$ QUBO $\rightarrow$ Ising.} \\
    Comprobar que la energía del Hamiltoniano de Ising construido reproduce exactamente la función objetivo financiera original (media-varianza con penalización de cardinalidad) para el conjunto completo de configuraciones binarias en instancias de tamaño reducido ($N \le 14$), empleando un solucionador exacto por fuerza bruta como referencia de verificación. \\
    \textit{Criterio de éxito:} Discrepancia numérica nula (dentro de la tolerancia de precisión de punto flotante de doble precisión, $\varepsilon \le 10^{-9}$) entre la energía de Ising y el valor de la función de coste QUBO en el 100\% de las configuraciones evaluadas.

    \item \textbf{OE2. Cuantificar la mejora en la factibilidad estructural de las soluciones cuánticas.} \\
    Determinar si la sustitución del mezclador transversal ($X$-Mixer) y la penalización cuadrática por un operador de intercambio $XY$ inicializado en un estado de Dicke incrementa de forma sustancial la proporción de muestras que satisfacen la restricción $\sum_i x_i = K$, frente al QAOA estándar, para un conjunto de instancias con $N \in [6, 20]$ y profundidades $p$ acotadas. \\
    \textit{Criterio de éxito:} Diferencia estadísticamente relevante y creciente con el tamaño del universo $N$ entre la tasa de factibilidad (\textit{Feasibility Ratio}) de ambos enfoques, verificable mediante el muestreo directo del vector de estado simulado.

    \item \textbf{OE3. Evaluar el efecto estabilizador de la inicialización TQA y la regularización Ridge $L_2$ sobre la convergencia clásica.} \\
    Medir si la combinación de una inicialización basada en \textit{Trotterized Quantum Annealing} (TQA) y una penalización continua $L_2$ sobre la trayectoria de los parámetros $(\gamma, \beta)$ reduce la variabilidad del resultado final del optimizador clásico (COBYLA) y/o el número de evaluaciones de la función de coste necesarias para alcanzar un óptimo local estable, en comparación con una inicialización aleatoria sin regularizar, sobre una instancia de referencia fija. \\
    \textit{Criterio de éxito:} Reducción observable y cuantitativa de la desviación estándar del \textit{Optimization GAP} entre distintas semillas aleatorias al introducir el esquema de regularización.

    \item \textbf{OE4. Contrastar el rendimiento del XY-QAOA regularizado frente a solucionadores de referencia clásicos.} \\
    Comparar el \textit{Optimization GAP} respecto al óptimo exacto determinista (Gurobi), el Ratio de Sharpe fuera de muestra (\textit{Out-of-Sample}) y el coste temporal de ejecución del XY-QAOA regularizado frente a Gurobi Optimizer y Simulated Annealing, utilizando series históricas reales de mercado, para determinar el rango de tamaños de universo $N$ en el que la propuesta cuántica resulta competitiva bajo simulación clásica exacta. \\
    \textit{Criterio de éxito:} Identificación explícita del umbral dimensional de $N$ a partir del cual el coste de la simulación de vectores de estado deja de ser computacionalmente práctico en CPU, junto con la caracterización cuantitativa de la relación entre precisión matemática (GAP) y robustez financiera en el régimen dimensional accesible.
\end{itemize}

\subsection{Hipótesis y Criterios de Éxito}
\label{subsec:hipotesis_criterios_exito}

Las hipótesis científicas derivadas de los objetivos específicos establecen las condiciones bajo las cuales se considerará demostrada la eficacia de la propuesta. Estos criterios serán contrastados empíricamente en el Capítulo~\ref{cap:evaluacion_resultados} a partir de las ejecuciones experimentales sobre datos de mercado, y su veredicto formal y cuantitativo de cumplimiento se recogerá de manera estructurada en el Capítulo~\ref{cap:conclusiones_lineas_futuras}. 

La Tabla~\ref{tab:hipotesis_exito} presenta la matriz de trazabilidad entre cada objetivo específico, su entregable o indicador de control asociado y la métrica de éxito predefinida.

\begin{table}[htbp]
    \centering
    \small
    \caption{Matriz de trazabilidad de objetivos específicos e hipótesis de éxito.}
    \label{tab:hipotesis_exito}
    \renewcommand{\arraystretch}{1.25}
    \begin{tabular}{cp{5.2cm}p{6.8cm}}
        \toprule
        \textbf{Objetivo} & \textbf{Entregable / Indicador} & \textbf{Métrica de Éxito} \\
        \midrule
        \textbf{OE1} & Verificación de equivalencia algebraica Markowitz $\rightarrow$ QUBO $\rightarrow$ Ising. & \textit{[PENDIENTE: consolidar tras Cap.~5 y 6 --- e.g., coincidencia analítica y numérica exacta ($\Delta < 10^{-9}$) en la totalidad de configuraciones del espectro evaluado].} \\
        \addlinespace
        \textbf{OE2} & Tasa de viabilidad matemática comparada: QAOA estándar vs. XY-QAOA. & \textit{[PENDIENTE: consolidar tras Cap.~6 con las mediciones empíricas --- e.g., preservación del 100\% de factibilidad frente al decaimiento monótono del enfoque penalizado].} \\
        \addlinespace
        \textbf{OE3} & Dispersión y estabilidad de la optimización clásica con/sin regularización Ridge $L_2$. & \textit{[PENDIENTE: consolidar tras Cap.~6 con la reducción de varianza y llamadas a función registradas entre semillas aleatorias].} \\
        \addlinespace
        \textbf{OE4} & Benchmarking multidimensional frente a Gurobi y Simulated Annealing (GAP, Sharpe OOS, Tiempo). & \textit{[PENDIENTE: consolidar tras Cap.~6, delimitando el rango dimensional donde la solución ofrece equilibrio entre robustez financiera y coste computacional].} \\
        \bottomrule
    \end{tabular}
\end{table}

% ==============================================================================
% 2.2. ALCANCE Y DELIMITACIÓN DEL ESTUDIO
% ==============================================================================
\section{Alcance y Delimitación del Estudio}
\label{sec:alcance_delimitacion}

Con el fin de asegurar la reproducibilidad metodológica y encuadrar con rigor científico el aporte del trabajo, se delimitan las inclusiones técnicas consideradas y se declaran explícitamente las exclusiones derivadas de las restricciones de la era NISQ.

\subsection{Fronteras del Proyecto e Inclusiones Técnicas}
\label{subsec:inclusiones_tecnicas}

El alcance del marco analítico y computacional desarrollado en esta memoria comprende las siguientes tres dimensiones operativas:

\begin{enumerate}
    \item \textbf{Modelado financiero con datos históricos reales:} Ingesta, depuración y modelización de series temporales de precios de cotización diarios de renta variable procedentes de mercados organizados e índices de referencia internacionales (S\&P 500 e IBEX 35). El detalle de la taxonomía de activos, las ventanas de calibración/prueba y los procedimientos de preprocesado e hiperparametrización se describe extensamente en el Capítulo~\ref{cap:marco_experimental}.
    
    \item \textbf{Simulación cuántica exacta mediante vectores de estado hasta $N=20$ cúbits:} Emulación matemática rigurosa de las transformaciones unitarias del circuito cuántico sobre arquitecturas de procesamiento central clásico (CPU), sin introducir simplificaciones de truncamiento de rango. La delimitación estricta en $N=20$ cúbits no constituye una elección metodológica arbitraria, sino que viene impuesta por la barrera física de memoria RAM ($2^N$ amplitudes complejas de doble precisión que en $N=20$ requieren almacenar vectores densos de más de un millón de componentes en cada evaluación del bucle variacional).
    
    \item \textbf{Soporte para restricciones duras de cardinalidad y presupuestarias:} Integración del confinamiento físico dentro del subespacio combinatorio mediante compuertas de intercambio XY para forzar la selección de exactamente $K$ activos, complementado con la penalización diagonal sobre el Hamiltoniano de coste para preservar el equilibrio presupuestario.
\end{enumerate}

\subsection{Exclusiones Metodológicas y Limitaciones NISQ}
\label{subsec:exclusiones_metodologicas}

Por imperativo de rigor frente a variables de ingeniería que exceden el diseño algorítmico del software, quedan formalmente excluidos los siguientes aspectos:

\begin{itemize}
    \item \textbf{Mitigación activa y corrección de errores en hardware físico real (QEC):} El diseño de códigos de corrección de errores activos (\textit{Fault-Tolerant Quantum Computing}) queda fuera del alcance práctico, dado que los dispositivos físicos actuales no disponen del número de cúbits físicos por cúbit lógico necesario para su implementación. La investigación se centra en la simulación exacta libre de ruido térmico, complementada con modelos analíticos de transpilación.
    
    \item \textbf{Optimización dinámica y rebalanceo multiperiodo:} El estudio se circunscribe a modelos estáticos de un único periodo (\textit{single-period asset allocation}). La programación dinámica estocástica y los costes de transacción por deslizamiento temporal en carteras continuas no forman parte del núcleo experimental.
    
    \item \textbf{Latencias de red y tiempos de encolamiento:} En la medición de tiempos computacionales se aísla el tiempo neto de cómputo del bucle híbrido CPU-simulador, descartando retardos de transmisión en red o esperas en planificadores de colas de plataformas de computación en la nube.
    
    \item \textbf{Estrategias de negociación de alta frecuencia (HFT):} El modelo no está orientado al arbitraje en milisegundos, sino a la asignación estructural y estratégica de activos bajo regímenes de rebalanceo periódico.
\end{itemize}

% ==============================================================================
% 2.3. ANTECEDENTES Y ESTADO DEL ARTE
% ==============================================================================
\section{Antecedentes y Estado del Arte}
\label{sec:antecedentes_estado_arte}

\subsection{De Markowitz al Problema Combinatorio}
\label{subsec:markowitz_combinatorio}

La teoría moderna de selección de carteras tiene su origen en el modelo de media-varianza propuesto por Harry Markowitz en 1952 \cite{Markowitz1952}. En su formulación original continua, el inversor busca un vector de ponderaciones que equilibre el retorno esperado frente al riesgo de la covarianza de los activos, resoluble de manera polinómica mediante algoritmos clásicos de programación cuadrática continua.

Sin embargo, cuando la operativa real de los mercados impone restricciones de cardinalidad discreta (seleccionar un número exacto $K$ de activos entre un universo de tamaño $N$), el problema requiere la incorporación de variables binarias de inclusión/exclusión. Esta modificación estructural altera radicalmente la complejidad analítica del modelo, transformándolo en un problema de Optimización Cuadrática Entera Mixta (MIQP), perteneciente a la clase de complejidad NP-hard. El tamaño del espacio de búsqueda discreto admisible queda gobernado por el coeficiente binomial $\binom{N}{K}$, cuya explosión combinatoria respecto al espacio completo $2^N$ fue tipificada cuantitativamente en la Tabla~\ref{tab:dimension_espacio} del Capítulo~\ref{cap:introduccion}. Ante este reto computacional, la industria financiera recurre tradicionalmente a dos familias de solucionadores clásicos.

\subsection{Solucionadores Clásicos de Referencia}
\label{subsec:solucionadores_clasicos}

En la práctica cuantitativa actual, la resolución de modelos MIQP de selección de carteras se fundamenta en dos aproximaciones complementarias:

\begin{itemize}
    \item \textbf{Solucionadores Deterministas Exactos (Gurobi Optimizer):} Arquitecturas de grado industrial que integran técnicas de ramificación y acotación (\textit{Branch-and-Bound}) con planos de corte poliédricos \cite{Gurobi2025}. Si bien garantizan la optimalidad global matemática y certifican el GAP exacto de la solución, su coste temporal en el peor de los casos escala de forma exponencial con el número de variables binarias, perdiendo viabilidad operativa cuando el universo de activos alcanza dimensiones elevadas.
    
    \item \textbf{Metaheurísticas Estocásticas (Simulated Annealing):} Algoritmos de búsqueda local inspirados en los procesos termodinámicos de recocido de materiales. Destacan por su rapidez temporal y su capacidad para abordar problemas combinatorios de gran escala, si bien carecen de certificados de convergencia global y presentan una inherente sensibilidad a la semilla pseudoaleatoria inicial y a los esquemas de enfriamiento.
\end{itemize}

Ambos paradigmas constituyen las referencias basales (\textit{baselines}) frente a las cuales se contrastará la metodología cuántica propuesta en los capítulos posteriores.

\subsection{Computación Cuántica Variacional: de Quantum Annealing a QAOA}
\label{subsec:qaoa_evolucion}

La incapacidad de los métodos clásicos deterministas para mitigar el crecimiento combinatorio en universos densos impulsó el interés por la computación cuántica. Históricamente, las primeras aproximaciones prácticas se desarrollaron sobre plataformas analógicas de recocido cuántico (\textit{Quantum Annealing}), diseñadas para minimizar funciones de energía mediante tunelamiento cuántico en redes de espines físicos. No obstante, las restricciones de conectividad en los grafos de acoplamiento físico y la susceptibilidad al ruido ambiental motivaron la transición hacia el paradigma cuántico digital basado en compuertas lógicas parametrizadas.

En este marco digital de la era NISQ \cite{Preskill2018}, el Algoritmo de Optimización Cuántica Aproximada (QAOA), introducido por Farhi et al. \cite{Farhi2014}, se posicionó como el referente fundamental para optimización combinatoria. QAOA opera alternando la aplicación de operadores unitarios parametrizados por ángulos $(\gamma, \beta)$ correspondientes a un Hamiltoniano de coste (que codifica la función objetivo) y un Hamiltoniano mezclador transversal ($X$-Mixer). En la formulación estándar, la observancia de restricciones duras como la cardinalidad $K$ se gestiona mediante términos de penalización cuadrática integrados en la función de coste, mapeando el problema a un formato de Optimización Binaria Cuadrática No Restringida (QUBO).

\subsection{Limitaciones Documentadas del QAOA Estándar}
\label{subsec:limitaciones_qaoa}

A pesar de su elegancia teórica, el QAOA estándar con penalizaciones presenta importantes limitaciones documentadas en la literatura especializada cuando se aplica a problemas fuertemente restringidos:

\begin{enumerate}
    \item \textbf{Deformación del paisaje energético por penalizaciones:} La necesidad de incorporar multiplicadores de penalización de gran magnitud para hacer cumplir la restricción $\sum_i x_i = K$ altera severamente la estructura del espectro de energía. La función de coste resultante exhibe profundos mínimos locales espurios y un paisaje altamente rugoso que dificulta la optimización clásica de los parámetros variacionales.
    
    \item \textbf{Mesetas estériles (\textit{Barren Plateaus}):} Como demostraron McClean et al. \cite{McClean2018}, los circuitos cuánticos parametrizados con inicializaciones aleatorias sufren un fenómeno por el cual la varianza del gradiente de la función de coste decae exponencialmente con el número de cúbits. Esto aplana el paisaje de optimización e impide que optimizadores locales clásicos guíen la convergencia del circuito.
\end{enumerate}

La confluencia de estas dos barreras computacionales motiva de manera directa la necesidad de abandonar el paradigma penalizado en favor de arquitecturas con preservación de simetría estructural.

% ==============================================================================
% 2.4. SÍNTESIS COMPARATIVA Y POSICIONAMIENTO DEL TRABAJO
% ==============================================================================
\section{Síntesis Comparativa y Posicionamiento del Trabajo}
\label{sec:sintesis_posicionamiento}

Frente a las limitaciones documentadas del QAOA estándar y la ineficiencia de evaluar el espacio completo de Hilbert $2^N$, el presente Trabajo de Fin de Máster adopta un enfoque de diseño algorítmico avanzado. En lugar de sancionar las soluciones no factibles mediante penalizaciones algebraicas externas, se restringe la dinámica física del circuito cuántico al subespacio geométrico estricto de las carteras válidas $\binom{N}{K}$, complementando dicha estructura con un esquema bicapa de inicialización adiabática guiada (TQA) y regularización Ridge $L_2$ en el optimizador clásico.

La Tabla~\ref{tab:estado_arte} sintetiza la comparativa taxonómica entre los diferentes paradigmas de optimización clásica y cuántica analizados en este estado del arte, posicionando conceptualmente la arquitectura propuesta dentro del panorama tecnológico contemporáneo.

\begin{table}[htbp]
    \centering
    \small
    \caption{Comparativa taxonómica del estado del arte en algoritmos de optimización combinatoria clásica y cuántica.}
    \label{tab:estado_arte}
    \renewcommand{\arraystretch}{1.3}
    \begin{tabular}{p{2.6cm}p{2.4cm}p{2.8cm}p{2.2cm}p{2.7cm}}
        \toprule
        \textbf{Algoritmo / Enfoque} & \textbf{Paradigma de Computación} & \textbf{Manejo de Restricciones ($K$)} & \textbf{Eficiencia en Cúbits / Recursos} & \textbf{Escalabilidad en la Era NISQ} \\
        \midrule
        \textbf{Gurobi Optimizer} \cite{Gurobi2025} & Clásico Determinista Exacto & Planos de corte y bifurcación ramificada. & Memoria RAM y CPU con escalado exponencial en peor caso. & Nula para instancias densas de gran escala ($N \gg 100$). \\
        \addlinespace
        \textbf{Simulated Annealing} & Clásico Estocástico Heurístico & Penalización o rechazo analítico de vecindarios. & Cómputo CPU secuencial ligero y altamente paralelizable. & Alta rapidez temporal, pero sin certificado de optimalidad global. \\
        \addlinespace
        \textbf{Standard QAOA} \cite{Farhi2014} & Cuántico Digital Variacional & Penalización cuadrática en la función de coste (QUBO). & Requiere $N$ cúbits sin registros auxiliares. & Baja (alta sensibilidad a mesetas estériles e inestabilidad por penalización). \\
        \addlinespace
        \textbf{Quantum Annealing} & Cuántico Adiabático Nativo & Mapeo de penalizaciones QUBO en grafos físicos. & Alta dependencia del grafo de acoplamiento del chip. & Media (condicionada por la conectividad limitada del hardware). \\
        \addlinespace
        \textbf{XY-QAOA Regularizado (Propuesto)} & Cuántico Híbrido Variacional & Confinamiento físico en subespacio combinatorio reducido (desarrollado en el Cap.~\ref{cap:marco_teorico}). & Optimizado a $N$ cúbits lineales libres de registros auxiliares. & Alta (menor profundidad de circuito y mayor resistencia al ruido físico \cite{Barkoutsos2020}). \\
        \bottomrule
    \end{tabular}
\end{table}

El Capítulo~\ref{cap:marco_teorico} retoma este posicionamiento taxonómico para formalizar matemáticamente la arquitectura XY-QAOA regularizada y derivar analíticamente los Hamiltonianos de coste y mezcla, cuya implementación algorítmica y validación experimental sobre series financieras reales se detallan en los Capítulos~\ref{cap:marco_experimental} y~\ref{cap:evaluacion_resultados}, respectivamente.
